% ----------------------
\subsection{Section 130 (2 April 1843)}
% ----------------------
	\label{DC130}
	
	% -----
	\subsubsection{The original source for DC 130}
	% -----
	
		Note that the JSP has two original sources for the remarks that make up section 130. The first, which seems to be the source the section actually comes from, is \hyperref[EMS-1843-JS-2-APR-1843-WC]{William Clayton's account}. The other source is \hyperref[EMS-1843-JS-2-APR-1843-WR]{Willard Richards' account}.
	
	% -----
	\subsubsection{Historical Background}
	% -----
		\label{DC130-HB}
		
		% --
		\paragraph{JSP}
		% --
		
			``On 2 April 1843, JS responded to questions and instructed the Saints in Macedonia, Illinois, about the nature of God, prophecies of the last days, and the afterlife. The previous day, he began a journey from Nauvoo, Illinois, to Macedonia and Carthage, Illinois. During this four-day trip, JS, accompanied by William Clayton, Orson Hyde, and Jacob B. Backenstos, preached to members of the branch of the church in Macedonia, visited JS's sister Sophronia Smith McCleary, and conducted business at Carthage. After spending the night in the Macedonia home of church members Benjamin F. and Melissa LeBaron Johnson, JS spent Sunday, 2 April 1843, attending various public gatherings and private meetings.

			``Following a Sunday morning service in which Orson Hyde preached on portions of the New Testament, JS corrected elements of Hyde's sermon. He also answered theological questions and related a dream he had, which Hyde interpreted. 
			
			``At some point `during the day,' JS made several `remarks on doctrine' relating to the timing of events prophesied to occur in the days leading up to the second coming of Jesus Christ. In the early 1840s, there was considerable interest in the prophecies relating to the latter days. Baptist minister William Miller, for example, predicted an approximate date for Christ's second coming based on a close reading of the Bible. He claimed that his examination of the books of Daniel and Revelation had revealed the second advent of Christ to be imminent, and his message was accepted by thousands of Christians in the early 1840s. Perhaps in response to a belief held by some Millerites that 3 April was the day appointed for the return of Christ, JS addressed this topic twice on 2 April.
			
			``Later that day, during Sabbath meetings held at 1:00 p.m. and again at 7:00 p.m., JS preached on the meanings of several verses from the book of Revelation. The subject may have been of interest because of a high council meeting held two weeks earlier, in which church member Pelatiah Brown was tried for teaching false doctrine because of his controversial interpretation of these same passages. 

			``During his evening discourse, JS again corrected Hyde's morning sermon. At the close of the discourse, he offered Hyde a moment to respond. The apostle stated that JS had said all that could be said and closed the meeting. JS spent the evening at the Johnsons' home, where he provided private instruction on portions of the book of Revelation relating to the sealing of the 144,000.
			
			``William Clayton, one of JS's travel companions, recorded details of JS's 2 April 1843 Sabbath-day teachings as he taught them. Clayton's journal contains other entries for which an earlier text is also extant, and comparisons of these entries indicate that when he copied his notes, he often expanded or otherwise changed them. Consequently, it seems likely that Clayton expanded the entries for 2 April 1843 as well. After the trip, Willard Richards, who kept JS's journal but was not present at the meetings, reconstructed a timeline for these teachings in JS's journal, presumably relying on Clayton's notes and conversations with people who were present. Because Clayton's initial notes are apparently not extant, it is impossible to know the exact sources Richards used to create his entry. Therefore, both Clayton's journal entry and the JS journal entry inscribed by Richards are featured here. Annotation that appears in Clayton's version of the discourse is not repeated in corresponding locations in the version by Richards. Richards included some details of the day's events that were not part of JS's teachings. They have been left to stand in the transcription.''			
					
	% -----
	\subsubsection{Versions}
	% -----
		\label{DC130-V}
		
		See the \href{https://drive.google.com/file/d/1goAvNz8jS4R8slYfMG_db55Ty2z_vmgK/view?usp=sharing}{sections comparison} document.
	
		\begin{compactdesc}
			\item[JSP] \hyperref[EMS-1843-JS-2-APR-1843]{2 April 1843}
			\item[BC] ---
			\item[DC 1835] ---
			\item[DC 1844] ---
			\item[DN] 6:18 (9 July 1856), 137
			\item[MS] 20:46 (13 November 1858), 727--728
		\end{compactdesc}

	% -----
	\subsubsection{Section 130:1} % *DC130-1 
	% -----
		\label{DC130-1}

		WHEN the Savior shall appear we shall see him as he is.%
			\footnote{%
				% \label{}%
				Compare \hyperref[MORONI-7-48]{Moroni 7:48}.
				}
		We shall see that he is a man like ourselves.

		% \begin{framed}
		% \scriptsize
			% \begin{compactdesc}
				% \item[JSP] 
				% % \item[BC] 
				% % \item[]
			% \end{compactdesc}
		% \end{framed}

	% -----
	\subsubsection{Section 130:2} % *DC130-2 
	% -----
		\label{DC130-2}

		And that same \hyperref[SOCIALITY]{sociality} which exists among us here will exist among us there, only it will be coupled with eternal glory, which glory we do not now enjoy.%
			\footnote{%
				\label{DC130-2-enjoy}%
				See JS's \hyperref[EMS-1836-JS-7-JAN-1836]{journal} for 7 January 1836---he goes to a feast and writes, ``before we parted we had some of the Songs of Zion sung, and our hearts were made glad while partaking of an antipast of those Joys that will be poured upon the head of the Saints w[h]en they are gathered together on Mount Zion to enjoy eachothers society forever more.'' By ``antipast'' he means ``antepast,'' which means ``a foretaste, and appetizer.''
				}

		% \begin{framed}
		% \scriptsize
			% \begin{compactdesc}
				% \item[JSP] 
				% % \item[BC] 
				% % \item[]
			% \end{compactdesc}
		% \end{framed}

	% -----
	\subsubsection{Section 130:3} % *DC130-3 
	% -----
		\label{DC130-3}

		\hyperref[JOHN-14-23]{John 14:23}---The appearing of the Father and the Son, in that verse, is a personal appearance; and the idea that the Father and the Son dwell in a man's heart is an old sectarian notion, and is false.

		% \begin{framed}
		% \scriptsize
			% \begin{compactdesc}
				% \item[JSP] 
				% % \item[BC] 
				% % \item[]
			% \end{compactdesc}
		% \end{framed}

	% -----
	\subsubsection{Section 130:4} % *DC130-4 
	% -----
		\label{DC130-4}

		In answer to the question---Is not the reckoning of God's time, angel's time, prophet's time, and man's time, according to the planet on which they reside?

		% \begin{framed}
		% \scriptsize
			% \begin{compactdesc}
				% \item[JSP] 
				% % \item[BC] 
				% % \item[]
			% \end{compactdesc}
		% \end{framed}

	% -----
	\subsubsection{Section 130:5} % *DC130-5 
	% -----
		\label{DC130-5}

		I answer, Yes. But there are no angels who minister to this earth but those who do belong or have belonged to it.

		% \begin{framed}
		% \scriptsize
			% \begin{compactdesc}
				% \item[JSP] 
				% % \item[BC] 
				% % \item[]
			% \end{compactdesc}
		% \end{framed}

	% -----
	\subsubsection{Section 130:6} % *DC130-6 
	% -----
		\label{DC130-6}

		The angels do not reside on a planet like this earth;

		% \begin{framed}
		% \scriptsize
			% \begin{compactdesc}
				% \item[JSP] 
				% % \item[BC] 
				% % \item[]
			% \end{compactdesc}
		% \end{framed}

	% -----
	\subsubsection{Section 130:7} % *DC130-7 
	% -----
		\label{DC130-7}

		But they reside in the presence of God, on a globe like a sea of glass and fire,%
			\footnote{%
				% \label{}%
				For this ``sea,'' see \hyperref[REVELATIONNT-15-2]{Revelation 15:2}.
				}
		where all things for their glory are manifest, past, present, and future, and are continually before the Lord.

		% \begin{framed}
		% \scriptsize
			% \begin{compactdesc}
				% \item[JSP] 
				% % \item[BC] 
				% % \item[]
			% \end{compactdesc}
		% \end{framed}

	% -----
	\subsubsection{Section 130:8} % *DC130-8 
	% -----
		\label{DC130-8}

		The place where God resides is a great Urim and Thummim.

		% \begin{framed}
		% \scriptsize
			% \begin{compactdesc}
				% \item[JSP] 
				% % \item[BC] 
				% % \item[]
			% \end{compactdesc}
		% \end{framed}

	% -----
	\subsubsection{Section 130:9} % *DC130-9 
	% -----
		\label{DC130-9}

		This earth, in its sanctified and immortal state, will be made like unto crystal and will be a Urim and Thummim to the inhabitants who dwell thereon,%
			\footnote{%
				% \label{}%
				See BY, \hyperref[EMS-1867-BY-27-AUG-1867]{27 August 1867}: ``we work away until we drive the power of Satan from the earth then the earth will be sanctified by the habitation and prayers of the saints angels Jesus [--?] day and [nor/wn?] night tending to the very business we are upon to purify and sanctify and prepare so that this earth may be redeemed and sanctified and by and by be glorious and become like a sea of glass so says John the Revelator and Joseph says be like a Urim and Thummin and those that live on it all to do look into the earth and see all things.''
				}
		whereby all things pertaining to an inferior kingdom, or all kingdoms of a lower order, will be manifest to those who dwell on it; and this earth will be Christ's.

		% \begin{framed}
		% \scriptsize
			% \begin{compactdesc}
				% \item[JSP] 
				% % \item[BC] 
				% % \item[]
			% \end{compactdesc}
		% \end{framed}

	% -----
	\subsubsection{Section 130:10} % *DC130-10 
	% -----
		\label{DC130-10}

		Then the white stone mentioned in \hyperref[REVELATIONNT-2-17]{Revelation 2:17}, will become a Urim and Thummim to each individual who receives one, whereby things pertaining to a higher order of kingdoms will be made known;

		% \begin{framed}
		% \scriptsize
			% \begin{compactdesc}
				% \item[JSP] 
				% % \item[BC] 
				% % \item[]
			% \end{compactdesc}
		% \end{framed}

	% -----
	\subsubsection{Section 130:11} % *DC130-11 
	% -----
		\label{DC130-11}

		And a white stone is given to each of those who come into the celestial kingdom, whereon is a new name written, which no man knoweth save he that receiveth it. The new name is the key word.

		% \begin{framed}
		% \scriptsize
			% \begin{compactdesc}
				% \item[JSP] 
				% % \item[BC] 
				% % \item[]
			% \end{compactdesc}
		% \end{framed}

	% -----
	\subsubsection{Section 130:12} % *DC130-12 
	% -----
		\label{DC130-12}

		I prophesy, in the name of the Lord God, that the commencement of the difficulties which will cause much bloodshed previous to the coming of the Son of Man will be in South Carolina.

		% \begin{framed}
		% \scriptsize
			% \begin{compactdesc}
				% \item[JSP] 
				% % \item[BC] 
				% % \item[]
			% \end{compactdesc}
		% \end{framed}

	% -----
	\subsubsection{Section 130:13} % *DC130-13 
	% -----
		\label{DC130-13}

		It may probably arise through the slave question. This a voice declared to me, while I was praying earnestly on the subject, December 25th, 1832.

		% \begin{framed}
		% \scriptsize
			% \begin{compactdesc}
				% \item[JSP] 
				% % \item[BC] 
				% % \item[]
			% \end{compactdesc}
		% \end{framed}

	% -----
	\subsubsection{Section 130:14} % *DC130-14 
	% -----
		\label{DC130-14}

		I was once praying very earnestly to know the time of the coming of the Son of Man, when I heard a voice repeat the following:

		% \begin{framed}
		% \scriptsize
			% \begin{compactdesc}
				% \item[JSP] 
				% % \item[BC] 
				% % \item[]
			% \end{compactdesc}
		% \end{framed}

	% -----
	\subsubsection{Section 130:15} % *DC130-15 
	% -----
		\label{DC130-15}

		Joseph, my son, if thou livest until thou art eighty-five years old, thou shalt see the face of the Son of Man; therefore let this suffice, and trouble me no more on this matter.

		% \begin{framed}
		% \scriptsize
			% \begin{compactdesc}
				% \item[JSP] 
				% % \item[BC] 
				% % \item[]
			% \end{compactdesc}
		% \end{framed}

	% -----
	\subsubsection{Section 130:16} % *DC130-16 
	% -----
		\label{DC130-16}

		I was left thus, without being able to decide whether this coming referred to the beginning of the millennium or to some previous appearing, or whether I should die and thus see his face.

		% \begin{framed}
		% \scriptsize
			% \begin{compactdesc}
				% \item[JSP] 
				% % \item[BC] 
				% % \item[]
			% \end{compactdesc}
		% \end{framed}

	% -----
	\subsubsection{Section 130:17} % *DC130-17 
	% -----
		\label{DC130-17}

		I believe the coming of the Son of Man will not be any sooner than that time.

		% \begin{framed}
		% \scriptsize
			% \begin{compactdesc}
				% \item[JSP] 
				% % \item[BC] 
				% % \item[]
			% \end{compactdesc}
		% \end{framed}

	% -----
	\subsubsection{Section 130:18} % *DC130-18 
	% -----
		\label{DC130-18}

		Whatever principle%
			\footnote{%
				% \label{}%
				Note that \hyperref[EMS-1843-JS-2-APR-1843-WC]{William Clayton's record} of this discussion uses the word ``principal,'' not ``principle.'' The other account, from \hyperref[EMS-1843-JS-2-APR-1843-WR]{Willard Richards}, has ``principle.'' ``Principal'' does appear in the JS concordance, but only around 10 times, and only once as a noun, that I can see (the other times it's an adjective).
				}
		of intelligence we attain unto in this life, it will rise with us in the resurrection.%
			\footnote{%
				% \label{}%
				The \hyperref[EMS-1843-JS-2-APR-1843-WR]{Willard Richards account} of this discussion has ``revalatin [revelation]'' here, instead of ``resurrection.'' The other account has ``resurrection.''
				}

		% \begin{framed}
		% \scriptsize
			% \begin{compactdesc}
				% \item[JSP] 
				% % \item[BC] 
				% % \item[]
			% \end{compactdesc}
		% \end{framed}

	% -----
	\subsubsection{Section 130:19} % *DC130-19 
	% -----
		\label{DC130-19}

		And if a person gains more knowledge and intelligence in this life through his diligence and obedience than another, he will have so much the advantage in the world to come.

		% \begin{framed}
		% \scriptsize
			% \begin{compactdesc}
				% \item[JSP] 
				% % \item[BC] 
				% % \item[]
			% \end{compactdesc}
		% \end{framed}

	% -----
	\subsubsection{Section 130:20} % *DC130-20 
	% -----
		\label{DC130-20}

			\footnote{%
				% \label{}%
				Compare the principle of these two verses with the one at \hyperref[DC132-5]{DC 132:5}.
				}%
		There is a law, irrevocably decreed in heaven before the foundations of this world, upon which all blessings are predicated---

		% \begin{framed}
		% \scriptsize
			% \begin{compactdesc}
				% \item[JSP] 
				% % \item[BC] 
				% % \item[]
			% \end{compactdesc}
		% \end{framed}

	% -----
	\subsubsection{Section 130:21} % *DC130-21 
	% -----
		\label{DC130-21}

		And when we obtain any blessing from God, it is by obedience to that law upon which it is predicated.

		% \begin{framed}
		% \scriptsize
			% \begin{compactdesc}
				% \item[JSP] 
				% % \item[BC] 
				% % \item[]
			% \end{compactdesc}
		% \end{framed}

	% -----
	\subsubsection{Section 130:22} % *DC130-22 
	% -----
		\label{DC130-22}

		The Father has a body of flesh and bones as tangible as man's; the Son also; but the Holy Ghost has not a body of flesh and bones, but is a personage of Spirit. Were it not so, the Holy Ghost could not dwell in us.%
			\footnote{%
				% \label{}%
				The verse as it stands here says almost the opposite from what the two original sources say. The \hyperlink{EMS-1843-JS-2-APR-1843-WC-l}{William Clayton account} (passage l) says, ``The Holy Ghost is a personage, and a person cannot have the personage of the H. G. in his heart. A man may receive the gifts of the H. G, and the H. G. may descend upon a man but not to tarry with him.'' The \hyperlink{EMS-1843-JS-2-APR-1843-WR-n}{Willard Richards account} (passage n) says, ``the Father has a body of flesh \& bones as tangible as mans the Son also, but the Holy Ghost is a personage of spirit.--- and a person cannot have the personage <of the H G. [Holy Ghost]> in his heart he may recive the gift of the holy Ghost. it may descend upon him but not to tarry with him.---''
				}

		% \begin{framed}
		% \scriptsize
			% \begin{compactdesc}
				% \item[JSP] 
				% % \item[BC] 
				% % \item[]
			% \end{compactdesc}
		% \end{framed}

	% -----
	\subsubsection{Section 130:23} % *DC130-23 
	% -----
		\label{DC130-23}

			\footnote{%
				% \label{}%
				See the note to the previous verse for more about the remark here.
				}%
		A man may receive the Holy Ghost, and it may descend upon him and not tarry with him.

		% \begin{framed}
		% \scriptsize
			% \begin{compactdesc}
				% \item[JSP] 
				% % \item[BC] 
				% % \item[]
			% \end{compactdesc}
		% \end{framed}